% Created 2026-06-06 Sat 16:26
% Intended LaTeX compiler: pdflatex
\documentclass[11pt]{article}
\usepackage[utf8]{inputenc}
\usepackage[T1]{fontenc}
\usepackage{graphicx}
\usepackage{longtable}
\usepackage{wrapfig}
\usepackage{rotating}
\usepackage[normalem]{ulem}
\usepackage{amsmath}
\usepackage{amssymb}
\usepackage{capt-of}
\usepackage{hyperref}
\author{Santiago Javier Proaño Yépez}
\date{\textit{<2026-06-06 Sat>}}
\title{Voltaje\textsubscript{en}\textsubscript{corriente}\textsubscript{alterna}}
\hypersetup{
 pdfauthor={Santiago Javier Proaño Yépez},
 pdftitle={Voltaje\textsubscript{en}\textsubscript{corriente}\textsubscript{alterna}},
 pdfkeywords={},
 pdfsubject={},
 pdfcreator={Emacs 29.3 (Org mode 9.6.15)}, 
 pdflang={English}}
\begin{document}

\maketitle
\tableofcontents

\section{Voltaje en corriente alterna}
\label{sec:org511c8a9}

Hay varias formas de representar el voltaje en una corriente sinusoidal.

\subsection{Voltaje instantáneo}
\label{sec:org5f33480}

Es el voltaje en cualquier momento.

[
\begin{align}
V_{t1} &= V_1 \
V_{t2} &= V_2
\end{align}
]

\subsection{Voltaje pico}
\label{sec:org807d260}

Es el voltaje más alto de una función.

\subsection{Voltaje pico a pico}
\label{sec:org40f8ac9}

Es la suma del pico positivo y el pico negativo.

[
V\textsubscript{pp}=V\textsubscript{pinf}+V\textsubscript{psup}=2V\textsubscript{p}
]

\subsection{Voltaje RMS (Root Mean Square)}
\label{sec:org2e4f976}

[
V\textsubscript{RMS}=\frac{V_p}\{\sqrt{2}\}
]

\subsection{Voltaje promedio}
\label{sec:orgc0afef9}

[
V\textsubscript{prom}=\frac{2}{\pi}V\textsubscript{p}
]
\end{document}
